\section{Apéndice: terminología y abreviaturas}
\subsection{Explicación de los términos principales}
\begin{table}[H]
\centering
\begin{tabular}{p{4cm}p{11cm}}
\toprule
Término & Explicación \\
\midrule
Chain-of-Thought & Hace que el modelo enumere los pasos de razonamiento intermedios, lo cual proporciona la base para el self-check. \\
Self-Rewarding & El modelo genera tareas por sí mismo y utiliza el reward para mejorar el reasoning, sin depender de labels manuales. \\
Self-Instruct & A partir de unos cuantos seeds, genera automáticamente instruction→input→output, con lo que amplía el corpus de entrenamiento. \\
Verified Response & La respuesta final que el modelo genera después de autoverificarse, la cual sirve como señal confiable para el reward. \\
Meta Judge & Un evaluator que compara el draft con el verified y decide cuál es más confiable. \\
SelfT Evaluators & Un conjunto de datos de judge entrenado usando el verified response como etiqueta. \\
\bottomrule
\end{tabular}
\caption{Varios términos mencionados repetidamente en las notas de esta clase.}
\end{table}

\subsection{Abreviaturas relacionadas}
\begin{itemize}
    \item PRM (Process Reward Model): recompensa paso a paso, que se usa para incentivar cada paso del reasoning.
    \item ORM (Outcome Reward Model): solo considera la respuesta final; es el enfoque típico del Meta judge.
    \item RLHF (Reinforcement Learning from Human Feedback): entrenamiento tradicional de reward que incluye retroalimentación humana.
    \item DPO (Direct Preference Optimization): una estrategia de actualización de parámetros que no usa RL y también aplica preference learning.
\end{itemize}

\subsection{Resumen de la sección}
El apéndice de terminología y abreviaturas ayuda a repasar las palabras clave de la clase, y facilita ubicar rápidamente la relación entre Self-Rewarding, Meta Judge y Self-Instruct.

